\documentclass[UTF8,fontset=fandol,11pt,a4paper]{ctexart}
\usepackage[a4paper,margin=1.75cm,headheight=15pt]{geometry}
\usepackage{amsmath,amssymb,mathtools,booktabs,tabularx,array,enumitem}
\usepackage[most]{tcolorbox}
\usepackage{tikz}
\usetikzlibrary{arrows.meta,positioning,calc,fit}
\usepackage{xcolor,fancyhdr,hyperref,microtype,lastpage}
\newcolumntype{Y}{>{\raggedright\arraybackslash}X}
\newcolumntype{L}[1]{>{\raggedright\arraybackslash}p{#1}}
\newcolumntype{C}[1]{>{\centering\arraybackslash}p{#1}}
\setlength{\parindent}{2em}
\setlength{\parskip}{0.34em}
\setlength{\emergencystretch}{3em}
\renewcommand{\arraystretch}{1.2}
\setlength{\tabcolsep}{3pt}
\setlist[itemize]{itemsep=0.18em,topsep=0.35em,leftmargin=2em}
\setlist[enumerate]{itemsep=0.18em,topsep=0.35em,leftmargin=2em}
\definecolor{deep}{RGB}{38,58,72}
\definecolor{soft}{RGB}{244,247,248}
\definecolor{greenish}{RGB}{52,91,74}
\definecolor{warnc}{RGB}{136,61,58}
\definecolor{purple}{RGB}{84,72,112}
\definecolor{rulegray}{RGB}{110,116,120}
\hypersetup{colorlinks=true,linkcolor=deep,urlcolor=purple,pdftitle={CO-08 总线与 I/O 硬件}}
\pagestyle{fancy}
\fancyhf{}
\lhead{\small\color{deep}\textbf{408 · 计算机组成原理 CO-08}}
\rhead{\small\color{deep}Request · Transaction · Completion}
\cfoot{\small 第 \thepage\ 页 / 共 \pageref{LastPage} 页}
\renewcommand{\headrulewidth}{0.4pt}
\newtcolorbox{corebox}[1]{breakable,colback=soft,colframe=deep,title=\textbf{#1},boxrule=0.8pt,arc=2mm}
\newtcolorbox{methodbox}[1]{breakable,colback=white,colframe=greenish,title=\textbf{#1},boxrule=0.7pt,arc=1.5mm}
\newtcolorbox{warnbox}[1]{breakable,colback=white,colframe=warnc,title=\textbf{#1},boxrule=0.7pt,arc=1.5mm}
\newcommand{\kw}[1]{\textbf{\color{deep}#1}}
\newcommand{\code}[1]{\texttt{#1}}
\title{\vspace{-1.1cm}\textbf{总线与 I/O 硬件}\\[0.4em]\large 同步处理器怎样与异步设备完成可确认的协作}
\author{408 · 计算机组成原理 Topic CO-08}
\date{}

\begin{document}
\maketitle

\begin{corebox}{母问题：一次设备工作怎样从意图变成可确认的完成？}
总线、控制器、轮询、中断和 DMA 不是五个孤立名词，而是一条责任链：
\[
\boxed{\text{Request}\to\text{Interface State}\to\text{Arbitration}\to\text{Transaction}\to\text{Transfer}\to\text{Completion Signal}}
\]
每一步都要回答：谁拥有控制权，谁驱动共享资源，数据经过哪里，什么状态证明本步结束。
\end{corebox}

\begin{methodbox}{本册学习范围}
\begin{itemize}
\item 本册解释多个主设备如何共享总线，设备控制器怎样接受请求、搬运数据并报告完成，以及查询、中断和 DMA 的硬件差异。
\item 本册重点解释事务四阶段、同步/异步握手、集中式/分布式仲裁、I/O 端口、中断入口、DMA 传送模式、通道和有效带宽。
\item 驱动程序、请求队列、阻塞/唤醒、缓冲区管理和设备调度属于软件；本册只说明硬件交给软件的状态和完成证据。
\end{itemize}
\end{methodbox}

\section{术语先行：总线和 I/O 的角色先分清}
\begin{methodbox}{本册的十一个关键词}
\begin{itemize}
\item \kw{总线：}连接多个主设备和目标设备、承载地址、数据与控制信号的共享或点对点互连。
\item \kw{主设备/目标设备：}当前发起事务的一方和被寻址、响应事务的一方；同一设备在不同事务中可能换角色。
\item \kw{事务：}从请求或授权开始，到数据传送和结束确认的一次完整总线操作。
\item \kw{仲裁：}多个主设备同时请求总线时，决定谁暂时获得驱动权的过程。
\item \kw{握手：}用 request/acknowledge、ready/wait 等信号确认双方何时发送和采样。
\item \kw{I/O 控制器：}夹在 CPU/总线和外设之间的硬件，提供数据、状态、命令寄存器并吸收速度差。
\item \kw{轮询 polling：}CPU 反复读取状态寄存器，直到设备报告就绪。
\item \kw{中断 interrupt：}设备提出请求，CPU 在允许的指令边界保存断点并转入中断服务程序。
\item \kw{DMA：}直接存储器访问；DMA 控制器代表设备取得总线，把一块数据直接搬到主存或从主存搬出。
\item \kw{DREQ/IRQ：}DREQ 是 DMA 传送请求，IRQ 是要求 CPU 处理事件的中断请求；两者都不是“已经完成”的证明。
\item \kw{MMIO：}把设备寄存器放进内存地址空间；访问是否可缓存、可重排和可用的宽度由平台属性规定。
\end{itemize}
\end{methodbox}

\tableofcontents
\newpage

\section{总线是协议，不是一捆静态导线}
共享互连至少包含 initiator/master、target/slave、arbiter 和 interconnect。一次事务通常经过请求、仲裁、地址/命令、一个或多个 data beat、response 与释放。DMA 控制器获授权后也可成为 initiator；“CPU 永远是主设备”不是不变量。

地址、数据与控制信息承担不同责任：地址标识目标，数据线宽度限制每 beat 原始数据量，控制/响应表达 read/write、byte enable、ready、error、request/grant 等协议状态。地址线分时复用不会自动增加数据带宽。

\section{同步、异步与分离事务}
同步总线按共同 clock 采样，慢目标可插入 wait state；异步总线靠 request/acknowledge 握手建立因果完成关系。二者描述\kw{定时协议}，串行/并行描述\kw{物理传输}，不能互推快慢。

分离事务把 request 与 response 解耦，长延迟时可释放通道服务其他请求，但需要事务 ID、缓冲和匹配。是否允许多个 outstanding request 是协议属性，不由“异步”一词推出。

\section{仲裁：选择共享资源的临时控制者}
仲裁只决定哪个请求者在当前时段获得发起权，不决定 OS 层的设备归属。固定优先级响应关键设备快，却可能饥饿；轮转改善公平性，却增加状态与最坏等待；链式、独立请求或分布式实现还要比较线路成本、仲裁延迟和故障范围。

\begin{center}\small
\begin{tabularx}{\linewidth}{L{2.2cm}YYY}\toprule
机制 & 关键状态 & 优势 & 风险\\\midrule
固定优先级 & pending 与 priority & 关键请求低延迟 & 低优先级饥饿\\
轮转 & last grant / pointer & 长期公平 & 状态与最坏等待增加\\
链式查询 & grant 传播位置 & 线路少 & 延迟累积、物理顺序固化\\
独立请求 & 每主设备 request/grant & 并行判优快 & 线路和逻辑成本高\\
\bottomrule
\end{tabularx}
\end{center}

\section{带宽题先写完整事务}
理论峰值为
\[
BW_{peak}=\frac{bytes}{beat}\times\frac{beats}{cycle}\times frequency.
\]
有效带宽还要扣除仲裁、地址/响应阶段、wait、方向切换、编码与空闲。若一个事务搬运 $D$ 字节、总耗时 $T$，稳定口径是 $BW_{effective}=D/T$。频率、总线周期、数据周期与设备时间必须分开列。

\begin{warnbox}{最小反例}
64 位、1 GHz 只给出每 beat 上限与时钟条件；若每次数据前有地址阶段、目标插入等待，或并非每周期都有 beat，有效带宽就不是自动的 8 GB/s。
\end{warnbox}

\section{I/O 接口把设备变成寄存器状态机}
控制器通常暴露 data/FIFO、status、control/command，以及可选地址、计数、描述符和 interrupt mask。软件写 command 只是启动；busy、ready、done、error 才描述后续状态。

memory-mapped I/O 用普通地址空间和 load/store；isolated I/O 用独立端口空间和专用指令。差别在软件可见寻址契约，不等于设备性能。MMIO 是否可缓存、能否重排和访问宽度须服从 ISA/平台规则。

\section{三种控制方式：谁等待、谁搬数据、谁通知}
\begin{center}\small
\begin{tabularx}{\linewidth}{L{2.0cm}YYY}\toprule
方式 & CPU 等待 & 数据路径 & 完成发现\\\midrule
Polling / PIO & 循环查询状态 & device $\leftrightarrow$ CPU register $\leftrightarrow$ memory & 主动读取 done\\
Interrupt-driven & 可执行其他任务 & 通常仍由 CPU 指令搬运 & 设备异步请求入口\\
DMA & 只配置与收尾 & device $\leftrightarrow$ memory & 块完成/错误通知\\
\bottomrule
\end{tabularx}
\end{center}
中断消除的是持续 busy waiting，不自动消除 ISR 和搬运成本；DMA 下沉的是传输循环，不代表 CPU、Cache、内存和互连没有成本。

\begin{center}
\providecolor{themebg}{HTML}{FAFAF7}
\providecolor{themefg}{HTML}{111111}
\providecolor{themegray}{HTML}{666666}
\providecolor{themecurve}{HTML}{1D63B8}
\providecolor{themealert}{HTML}{C53030}
\providecolor{themeamber}{HTML}{B86E00}
\providecolor{themegreen}{HTML}{15803D}
\begin{tikzpicture}[font=\scriptsize,color=themefg,text=themefg,>=Stealth]
\tikzset{
  colbox/.style={draw=themegray!60,rounded corners=5pt,fill=themefg!2!themebg,inner sep=8pt},
  box/.style={draw=themegray,rounded corners=3pt,minimum width=3.4cm,minimum height=.78cm,align=center,fill=themebg},
  cpu/.style={box,draw=themealert,fill=themealert!8!themebg,font=\scriptsize\bfseries,text=themealert},
  dev/.style={box,draw=themecurve,fill=themecurve!8!themebg,font=\scriptsize\bfseries,text=themecurve},
  dma/.style={box,draw=themeamber,fill=themeamber!10!themebg,font=\scriptsize\bfseries,text=themeamber},
  mem/.style={box,draw=themegreen,fill=themegreen!8!themebg,font=\scriptsize\bfseries,text=themegreen},
  flow/.style={->,line width=1.1pt,themecurve},
  cpuflow/.style={->,line width=1.1pt,themealert},
  dmaflow/.style={->,line width=1.1pt,themeamber},
  memflow/.style={->,line width=1.1pt,themegreen},
  note/.style={rounded corners=4pt,draw=themegray!70,fill=themefg!3!themebg,inner xsep=8pt,inner ysep=7pt,align=left}
}

% ── 标题与例题说明 ──
\node[anchor=west,font=\bfseries\Large] at (0,8.6) {Polling、Interrupt、DMA：CPU 等待、搬运与通知责任的演进};
\node[anchor=west,text=themegray,text width=16.8cm,align=left] at (0,8.0)
  {三种方式都需要设备接口状态；演进核心是逐步解放 CPU，将“轮询等待”与“逐字搬运”从 CPU 卸载到专用硬件。};

% ══════════════════════════════════════════════════════════════
% ── 第 1 列：程序查询 (Polling / PIO) (x = 2.4) ──
% ══════════════════════════════════════════════════════════════
\node[font=\bfseries\normalsize,text=themealert] at (2.4,7.4) {1. 程序查询 (Polling)};
\node[cpu] (p1) at (2.4,6.4) {CPU 循环读取 Status\\(Busy-waiting 盲等)};
\node[dev] (p2) at (2.4,5.0) {I/O 接口 Device Ready};
\node[cpu] (p3) at (2.4,3.6) {CPU 执行 Load/Store\\逐字搬运数据};
\node[mem] (p4) at (2.4,2.2) {主存 Memory};

\draw[cpuflow] (p1) -- node[right=2pt,font=\tiny\bfseries,text=themealert] {轮询} (p2);
\draw[cpuflow] (p2) -- node[right=2pt,font=\tiny\bfseries,text=themealert] {就绪} (p3);
\draw[cpuflow] (p3) -- node[right=2pt,font=\tiny\bfseries,text=themealert] {写入} (p4);

% ══════════════════════════════════════════════════════════════
% ── 第 2 列：程序中断 (Interrupt-driven) (x = 8.0) ──
% ══════════════════════════════════════════════════════════════
\node[font=\bfseries\normalsize,text=themecurve] at (8.0,7.4) {2. 程序中断 (Interrupt)};
\node[cpu] (i1) at (8.0,6.4) {CPU 执行其他程序\\(无需轮询等待)};
\node[dev] (i2) at (8.0,5.0) {设备就绪 发送 IRQ 中断};
\node[cpu] (i3) at (8.0,3.6) {CPU 暂停并执行 ISR\\逐字搬运数据};
\node[mem] (i4) at (8.0,2.2) {主存 Memory};

\draw[flow] (i1) -- node[right=2pt,font=\tiny\bfseries,text=themecurve] {并发} (i2);
\draw[flow] (i2) -- node[right=2pt,font=\tiny\bfseries,text=themecurve] {触发 ISR} (i3);
\draw[cpuflow] (i3) -- node[right=2pt,font=\tiny\bfseries,text=themealert] {搬运} (i4);

% ══════════════════════════════════════════════════════════════
% ── 第 3 列：直接内存访问 (DMA) (x = 13.8) ──
% ══════════════════════════════════════════════════════════════
\node[font=\bfseries\normalsize,text=themeamber] at (13.8,7.4) {3. 直接内存访问 (DMA)};
\node[cpu] (d1) at (13.8,6.4) {CPU 初始化 DMAC\\(首地址 / 块大小 / 控制)};
\node[dma] (d2) at (13.8,5.0) {DMAC 申请总线并接管\\(周期挪用 / 停止 CPU)};

% 硬件直接成块传输
\node[box,draw=themeamber,fill=themeamber!6!themebg,minimum width=4.2cm] (d_data) at (13.8,3.6) {
  \begin{tabular}{c}
    \textbf{硬件直接成块搬运 (Burst/Cycle Steal)}\\
    $\text{Device}\;\longleftrightarrow\;\text{硬件数据流}\;\longleftrightarrow\;\text{Memory}$
  \end{tabular}
};

\node[cpu] (d3) at (13.8,2.2) {CPU 中断服务程序 (ISR)\\仅作传输完成收尾};

\draw[dmaflow] (d1) -- node[right=2pt,font=\tiny\bfseries,text=themeamber] {启动} (d2);
\draw[dmaflow] (d2) -- node[right=2pt,font=\tiny\bfseries,text=themeamber] {接管总线} (d_data);
\draw[dmaflow] (d_data) -- node[right=2pt,font=\tiny\bfseries,text=themeamber] {整块传输完毕} (d3);

% ── 底部总结与误读警示 ──
\node[note,anchor=west,text width=16.8cm] at (0,0.9)
  {\textbf{责任分工对比：}\\
   $\bullet$ \textbf{Polling}：CPU 承担“等待就绪”与“逐字搬运”全部开销（CPU 利用率极低）；\\
   $\bullet$ \textbf{Interrupt}：解除 CPU 盲等，但每次 I/O 仍由 CPU ISR 逐字搬运，高频中断会引起巨大的上下文切换开销；\\
   $\bullet$ \textbf{DMA}：整块数据由 DMAC 硬件在设备与内存间直传，CPU 仅介入“启动初始化”与“完成收尾”两次。};

\node[anchor=west,font=\scriptsize\bfseries,text=themeamber,text width=16.8cm,align=left] at (0,-0.2)
  {概念区分：DREQ（设备向 DMAC 请求）$\ne$ HRQ/Bus Grant（DMAC 向 CPU 申请总线）$\ne$ IRQ（DMAC 向 CPU 发送完成中断）。};
\end{tikzpicture}

\end{center}

\section{中断硬件路径与软件交接}
\[
\text{device event}\to\text{pending}\to\text{mask/priority}\to\text{precise boundary}\to\text{architectural entry state}\to\text{vector}.
\]
外部中断与当前指令异步；异常通常与当前指令同步相关。CPU 是否自动关中断、保存哪些最小状态、向量怎样形成均由 ISA 规定。完整通用寄存器保存、驱动服务、确认设备、唤醒进程和调度属于软件。

中断控制器的 pending、enable/mask、priority 和 in-service 状态必须分开：请求到达不等于立即接受，接受不等于 handler 已完成，handler 返回也不必自动清除设备源头。

\subsection{可屏蔽中断与硬件响应条件}
对可屏蔽外部中断，硬件在响应前必须同时满足三个硬性判据：
\begin{enumerate}
  \item \kw{请求提出}：中断源接口电路确实已将中断请求信号置位；
  \item \kw{允许响应}：CPU 当前处于开中断状态（PSW 中 \code{IF = 1}）；
  \item \kw{指令边界}：当前执行指令已到达允许响应的周期边界（通常为执行周期的最后时刻）。
\end{enumerate}
不可屏蔽中断（NMI）可绕过 IF 屏蔽位，但仍须在指令执行边界响应；同步异常则由当前指令执行直接触发。

\subsection{硬件自动入口 vs ISR 软件服务}
硬件与软件的交接边界必须绝对分清：
\begin{itemize}
  \item \kw{硬件中断隐指令 (Hardware Auto Actions)}：关中断 $\to$ 保存断点（PC 与 PSW） $\to$ 形成中断服务程序入口地址并送入 PC；
  \item \kw{软件 ISR 服务 (Software Driver Service)}：保存通用寄存器现场 $\to$ 开中断（若允许嵌套） $\to$ 执行具体设备 I/O 服务 $\to$ 关中断 $\to$ 恢复通用寄存器 $\to$ 执行特权返回指令 \code{iret}。
\end{itemize}
向量中断中，向量号索引向量表得到入口地址；非向量中断由软件轮询源头。注意：“向量地址”（表项地址）与“中断向量”（表项内容/入口地址）不可混为一谈。

\begin{center}
\providecolor{themebg}{HTML}{FAFAF7}
\providecolor{themefg}{HTML}{111111}
\providecolor{themegray}{HTML}{666666}
\providecolor{themecurve}{HTML}{1D63B8}
\providecolor{themealert}{HTML}{C53030}
\providecolor{themeamber}{HTML}{B86E00}
\providecolor{themegreen}{HTML}{15803D}
\begin{tikzpicture}[font=\scriptsize,color=themefg,text=themefg,>=Stealth]
\tikzset{
  hw/.style={draw=themecurve,rounded corners=3pt,minimum width=3.4cm,minimum height=.82cm,align=center,fill=themecurve!8!themebg,font=\scriptsize\bfseries,text=themecurve},
  sw/.style={draw=themeamber,rounded corners=3pt,minimum width=3.4cm,minimum height=.82cm,align=center,fill=themeamber!8!themebg,font=\scriptsize\bfseries,text=themeamber},
  dev/.style={draw=themegray,rounded corners=3pt,minimum width=3.4cm,minimum height=.82cm,align=center,fill=themefg!3!themebg,font=\scriptsize\bfseries},
  flow/.style={->,line width=1.1pt,themecurve},
  swflow/.style={->,line width=1.1pt,themeamber},
  lbl/.style={font=\tiny\bfseries,fill=themebg,inner sep=1.5pt,rounded corners=1pt},
  note/.style={rounded corners=4pt,draw=themegray!70,fill=themefg!3!themebg,inner xsep=8pt,inner ysep=7pt,align=left}
}

% ── 标题与例题说明 ──
\node[anchor=west,font=\bfseries\Large] at (0,8.6) {一次外部中断：硬件负责形成合法入口，软件负责服务并恢复};
\node[anchor=west,text=themegray,text width=17.5cm,align=left] at (0,8.0)
  {中断响应四阶段：中断请求 $\to$ 硬件中断隐指令（判优/关中断/保护断点/引出向量）$\to$ 软件 ISR 执行 $\to$ 中断返回。};

% ── 第一阶段：硬件所有边界 (Hardware-Owned Boundary) ──
\node[dev] (event) at (2.2,6.5) {设备事件触发\\IRQ Pending 置 1};
\node[hw] (elig) at (7.8,6.5) {硬件中断判优与屏蔽\\Mask / Priority / 指令边界检测};
\node[hw] (entry) at (13.4,6.5) {中断隐指令 (硬件自动完成)\\关中断 + 压栈 PC/PSW + 向量寻址};

\draw[flow] (event) -- node[above=2pt,lbl,text=themegray] {发中断请求} (elig);
\draw[flow] (elig) -- node[above=2pt,lbl,text=themecurve] {允许响应 (Accept)} (entry);

% ── 第二阶段：软件服务阶段 (Software-Owned Service) ──
\node[sw] (save) at (13.4,4.7) {ISR 前导 (软件保护现场)\\压栈通用寄存器 + 开中断};
\node[sw] (service) at (7.8,4.7) {设备驱动与 ISR 服务\\读取数据 / 清除设备中断标志};
\node[sw] (restore) at (2.2,4.7) {ISR 后续 (软件恢复现场)\\恢复通用寄存器 + 关中断};

\draw[swflow] (entry) -- node[right=2pt,lbl,text=themeamber] {跳转 ISR} (save);
\draw[swflow] (save) -- node[above=2pt,lbl,text=themeamber] {执行服务} (service);
\draw[swflow] (service) -- node[above=2pt,lbl,text=themeamber] {服务完毕} (restore);

% ── 第三阶段：硬件特权返回与恢复执行 ──
\node[hw] (iret) at (7.8,2.9) {中断返回指令 (IRET)\\硬件自动弹出 PC / PSW 并恢复特权级};
\node[dev] (resume) at (13.4,2.9) {恢复原被中断程序继续执行};

\draw[swflow] (restore.south) to[out=-90,in=180] node[below=2pt,lbl,text=themeamber] {执行 IRET} (iret.west);
\draw[flow] (iret) -- node[above=2pt,lbl,text=themecurve] {恢复控制流} (resume);

% 虚线包络区域标注
\draw[themecurve,dashed,rounded corners=6pt,line width=0.8pt] (0.3,5.8) rectangle (15.3,7.2);
\node[font=\scriptsize\bfseries,text=themecurve,anchor=west] at (0.5,7.0) {硬件中断响应阶段 (Hardware)};

\draw[themeamber,dashed,rounded corners=6pt,line width=0.8pt] (0.3,4.0) rectangle (15.3,5.4);
\node[font=\scriptsize\bfseries,text=themeamber,anchor=west] at (0.5,5.2) {软件 ISR 服务阶段 (Software)};

% ── 底部总结与误读警示 ──
\node[note,anchor=west,text width=17.5cm] at (0,1.3)
  {\textbf{硬件与软件职责边界：}\\
   1. \textbf{硬件自动完成（中断隐指令）}：关中断 $\to$ 保存断点（PC 与 PSW）$\to$ 引出中断服务程序入口（向量地址）；\\
   2. \textbf{软件自主完成（ISR 代码）}：保存通用寄存器现场 $\to$ 实际设备 I/O 处理 $\to$ 清除中断源标志 $\to$ 恢复寄存器现场 $\to$ 执行 \texttt{IRET} 指令；\\
   3. \textbf{IRET 硬件恢复}：\texttt{IRET} 指令由硬件自动弹出栈中保存的 PC 与 PSW，原子恢复中断前的程序运行状态。};

\node[anchor=west,font=\scriptsize\bfseries,text=themealert,text width=17.5cm,align=left] at (0,0.15)
  {禁止误读：IRQ pending=1 仅表示设备发起了请求，不等于 CPU 立即进入 ISR；屏蔽字、优先级以及当前是否处于不可打断指令周期都会影响响应时刻。};
\end{tikzpicture}

\end{center}

\subsection{响应优先级 vs 处理优先级}
\begin{itemize}
  \item \kw{响应优先级}：由硬件排队器或仲裁电路决定，不可软件动态更改；
  \item \kw{处理优先级}：由中断屏蔽字（Mask）在 ISR 中动态设定，可改变多重中断时的服务抢占次序。
\end{itemize}
若允许嵌套，屏蔽字至少应屏蔽“处理优先级不高于当前级”的中断（含自身），位序必须按题面编号严格重排。

\begin{methodbox}{硬件边界}
本 Topic 的最后一个稳定动作是“按 ISA 规定形成可进入软件的体系结构状态”。“唤醒哪个任务、是否重试请求、怎样处理错误”不写成硬件动作。
\end{methodbox}

\section{DMA：把搬运循环下沉到控制器}
CPU 先配置 source、destination、length、direction/control 或 descriptor；DMAC 获得设备和总线条件后发起 transaction，逐 beat 更新地址与计数，完成或错误时置状态并通知 CPU：
\[
\text{Setup}\to\text{DREQ}\to\text{Arbitrate}\to\text{Transfer}\to\text{Count/Address Update}\to\text{Complete}.
\]

burst 一次授权连续传输，吞吐高但占用长；cycle stealing 每次占少量周期，在 CPU 干扰与吞吐间折中；transparent/interleaved 利用约定时隙，依赖具体时序。DREQ 是设备到 DMAC 的传输请求，不是 CPU 中断；总线 grant 也不是上下文切换。

scatter-gather 可由描述符链减少大块非连续传输的重复配置，但描述符建立、地址映射、缓存一致性和完成处理仍产生 CPU/系统成本，因此“DMA 对任意传输恒为 $O(1)$ CPU 开销”不是稳定结论。

\section{贯穿母例：DMA 读取一个块}
\begin{enumerate}
\item 软件准备设备可用的目标地址、长度和方向，并按平台规则处理映射/一致性；
\item CPU 写设备/DMAC 控制寄存器，控制器进入 busy；
\item 设备产生 DREQ，DMAC 作为潜在 initiator 请求互连；
\item arbiter 授权后，DMAC 与内存/设备完成若干 beat，更新地址和剩余计数；
\item 共享内存或总线冲突表现为等待，CPU 是否同时运行取决于其当前资源需求；
\item 计数归零或错误使控制器写 completion state，并产生可屏蔽的 interrupt request；
\item CPU 在精确边界接受中断，形成 ISA 规定的入口状态；
\item 软件确认状态、解除映射、完成请求或唤醒任务。
\end{enumerate}

\section{五列概念边界}
\begin{center}\small
\begin{tabularx}{\linewidth}{L{1.8cm}C{0.35cm}L{1.9cm}YY}\toprule
概念 A & $\neq$ & 概念 B & 真正区别与题面信号 & 混淆后果\\\midrule
request & $\neq$ & transaction & 提出需求 vs 已获权并执行协议 & 漏仲裁与等待\\
bus grant & $\neq$ & interrupt accept & 互连所有权 vs CPU 控制流入口 & 把暂停访存写成 ISR\\
DREQ & $\neq$ & IRQ & 设备请求 DMAC 搬运 vs 控制器通知 CPU & 错画信号方向\\
interrupt entry & $\neq$ & OS service & ISA 最小入口 vs 软件处理与唤醒 & 硬件越界拥有 OS\\
peak bandwidth & $\neq$ & effective bandwidth & 理想 beat 速率 vs 完整事务有效数据率 & 漏协议开销\\
MMIO address & $\neq$ & memory data & 设备寄存器契约 vs 普通可缓存对象 & 错用缓存/重排\\
\bottomrule
\end{tabularx}
\end{center}

\section{做题控制协议}
\begin{enumerate}
\item 标 initiator、target、数据方向、共享资源与 completion 条件；
\item 把 request、arbitration、address/control、data beat、response 分开；
\item 逐阶段写资源控制者和允许重叠的 CPU/设备工作；
\item 带宽先算完整事务有效数据/总时间，再与峰值比较；
\item I/O 方式题回答谁等待、谁搬数据、谁通知；
\item 中断题分 pending、accept、entry、software service、return；
\item DMA 题分 setup、transfer、completion，并单列一致性和 OS handoff。
\end{enumerate}

\begin{methodbox}{压缩成第一动作}
看到总线或 I/O 题，先问：\kw{当前谁拥有控制权？当前事件只是 request，还是已经形成 transaction 或 completion evidence？}
\end{methodbox}

\section{本册与 CPU/设备的接口}
\begin{corebox}{从请求到完成证据}
本册负责共享线路、控制器状态、仲裁、事务和完成证据。中断解决“完成后怎样通知 CPU”，DMA 解决“数据怎样批量搬运”；两者可以组合但不能互相替代。请求队列、完成后的阻塞/唤醒和设备调度属于软件。
\end{corebox}
本册 Cost 要把事务有效数据、地址/控制/等待周期、仲裁、方向切换和可重叠阶段写全；总线峰值、有效带宽、设备 latency 与 CPU 指令时间不是同一个量。若访问是 MMIO，属性可能绕过普通 Cache，但仍需回到 CO-07/B02 的权限与地址接口。

\section{知识地图：从请求到完成证据}
\begin{center}\small
\begin{tabularx}{\linewidth}{L{3.0cm}YY}\toprule
考点 & 本册机制 & 接口\\\midrule
总线结构/仲裁/定时 & 角色、事务、握手、有效带宽 & CO-03/05\\
I/O 接口与编址 & 寄存器状态机、MMIO/isolated I/O & CO-02\\
轮询/中断 & 等待、搬运、通知；硬件精确入口 & X-B01/X-B03\\
DMA & setup、仲裁、transfer、completion & OS-05/X-B03\\
\bottomrule
\end{tabularx}
\end{center}

\begin{corebox}{一句话}
总线题追踪“谁在何时拥有共享线路”，I/O 题追踪“谁推进状态、谁搬数据、什么证据证明完成”；中断与 DMA 都只是跨层 handoff，不替 OS 完成软件责任。
\end{corebox}

本册把总线结构、I/O 控制方式、中断和 DMA 统一到“请求、授权、传输、完成”的状态机中。具体总线标准的电气参数、设备驱动策略和操作系统调度必须以题设或平台说明为准。

\section{补充细节：总线协议的可计算细节}
一次总线事务可按考纲的四阶段复述为：\kw{申请/分配}（主设备提出请求并由仲裁授权）、\kw{寻址}（发出从设备地址与命令）、\kw{传输}（一个或多个数据 beat）、\kw{结束}（撤销信号并释放总线）。在单主设备系统中第一、四阶段可以被硬件常态化隐藏，但 DMA、多处理器或多个 I/O 主设备出现后不能省略。突发传送只发首地址并自动递增后续地址，非突发传送则每个数据都重新付出寻址开销。

\begin{center}\small
\begin{tabularx}{\linewidth}{L{2.4cm}YY}
\toprule
信号组/对象 & 语义 & 易错边界\\\midrule
数据线 & 被搬运的数据、状态字或中断类型号等内容 & “命令字”是内容，仍走数据线；请求/应答是控制信号\\
地址线 & 目标存储单元或 I/O 端口编号 & 地址/数据复用减少引脚，不会自动提高数据阶段带宽\\
控制线 & 读写、时钟、BR/BG、ready/error、request/ack 等协议状态 & 控制线有出有入，不应按“全部单向”处理\\
\bottomrule
\end{tabularx}
\end{center}

总线还可以按距离分成片内、系统和通信/外部总线；不同速度等级之间由桥接器做协议与速率转换。总线宽度 $W$（字节）和工作传输频率 $F$ 决定理论峰值 $BW=W\times F$。若题目给的是基础时钟 $f$，且每周期有 $k$ 个数据传送，则 $F=kf$；若给的是 MT/s 或 GT/s 这类有效传输率，就不要再次乘 $k$。带宽只描述数据阶段的上界，地址、仲裁、等待和方向切换应放进另一个“完整事务有效速率”计算。

\section{四种定时协议与异步互锁}
同步通信让所有动作对齐公共时钟，短距离且部件速度相近时控制简单；慢设备可插入固定的 wait state，但固定时钟仍按约定的最慢时间运行。半同步保留时钟边沿，在从设备来不及时由 WAIT 拉长一个或多个整周期，适合速度有差异而仍希望保留同步实现的系统。

异步通信没有公共时钟，依靠 request/acknowledge 握手。按信号撤销是否等待对方，可用互锁处数判断三类：不互锁（双方都按约定延时撤销，最快但最不可靠）、半互锁（主设备等回答后撤销请求，从设备不等请求撤销）和全互锁（请求等回答，回答再等请求撤销，最可靠但握手最长）。若单程传播及判别时间为 $t$、从设备准备时间为 $T_s$，在题面把 $t$ 明确为单程时，典型关键路径可写成
\[
T_{\text{no-lock}}\approx t+T_s,\qquad
T_{\text{half-lock}}\approx 2t+T_s,\qquad
T_{\text{full-lock}}\approx 4t+T_s.
\]
不要把一次往返时间误当单程，也不要把 $T_s$ 按互锁次数重复计算。这里的“速度—可靠性”排序只属于握手协议，不是串行/并行传输的排序。

\section{仲裁实现的线数—延迟—公平性三角}
集中式仲裁的三类经典实现可以按授权信息的编码理解：链式查询把授权逐跳传递，控制线数近似常数但响应慢、优先级由物理顺序固定且容易因断链影响下游；计数器定时查询广播设备编号，控制线约为 $\lceil\log_2 n\rceil+2$，可从 0 开始实现固定优先级，也可从上次终点继续实现轮转；独立请求为每个主设备配置 BR/BG，线数约 $2n+1$，判优最快且最灵活。分布式仲裁则把仲裁逻辑分散在各主设备，通过共享仲裁号比较避免单点仲裁器，代价是每个设备都需额外逻辑。

\begin{methodbox}{读图判仲裁}
先找“谁产生 BR、谁产生 BG、谁决定赢家”。若授权沿设备链传递，是链式；若设备地址在公共线上逐个比较，是计数器查询；若每台设备有独立 BR/BG，是独立请求；若不存在中央决策点而各设备同时比较仲裁号，是分布式。仲裁只决定临时总线控制者，不决定设备驱动归属。
\end{methodbox}

\section{I/O 接口与端口状态机的完整边界}
外设与 CPU 存在速度、信号格式和控制协议三重不匹配，CPU 实际对话的对象是设备控制器而非设备本体。接口向 CPU 暴露的端口通常分为数据/FIFO、状态和命令/控制三类；DMA 还需要把主存地址、长度、方向或描述符交给控制器。无条件传送的设备可以在约定时刻直接传输；条件传送则必须先读 ready/busy 等状态，再读写数据端口。

程序查询方式中，CPU 反复读取状态位直到就绪，设备准备数据期间两者串行；定时查询的 CPU 占比可按
\[
\rho_{poll}=\frac{R/u\times C_{poll}}{f}
\]
估算，其中 $R$ 是字节率，$u$ 是每次端口传送的字节数，$C_{poll}$ 是一次查询消耗的周期数，$f$ 是主频。$u$ 来自端口/缓冲寄存器宽度，不能默认等于 1 字节。I/O 指令能直接改变设备状态，通常属于特权操作；MMIO 是否允许缓存、重排或宽度变化则必须服从 ISA/平台属性。

\section{中断的五阶段与 CPU 成本}
程序中断可拆成“请求、判优、响应、服务、返回”五阶段。接口在设备就绪时置请求触发器；中断控制器按 pending、mask/enable、priority 和 in-service 状态判优；CPU 在符合 ISA 的精确边界执行中断隐指令，保存断点、形成入口并交给 ISR；ISR 再保护现场、读取/写入端口、确认设备、恢复现场并执行返回指令。请求到达不等于已经接受，handler 返回也不一定自动清除设备源头。

若每次服务涉及 $m$ 个周期、设备每秒产生 $R/u$ 次传送，单次中断 CPU 占比为
\[
\rho_{irq}=\frac{(R/u)\times m}{f}.
\]
中断方式的价值是释放设备准备期间的等待时间，而不是保证单次搬运比查询便宜；当设备到达间隔小于一次中断响应加服务时间时，中断在物理上就无法及时消费数据，应升级到 DMA 或更强的缓冲结构。

\section{DMA 控制器与三种传送方式}
DMA 控制器的最小寄存器集合可从“独立完成一块访存”推出：AR 保存当前主存地址并按 DR 字节数递增，WC 保存剩余传送个数并递减，DR 作为设备侧与总线侧的过渡缓冲，CR 保存方向、启动和状态；还需要总线请求/授权逻辑和完成中断电路。完整流程仍是预处理（CPU 设置参数）、数据传送（DMAC 发地址与读写信号并更新 AR/WC）、后处理（完成/错误中断）。

CPU 与 DMA 争用主存时有三种经典方式：\kw{停止 CPU 访存}控制简单，适合高速设备的成组传送但 CPU 整段停顿；\kw{周期挪用}每次只占一个或几个存储周期，设备两个数据之间的间隔必须足以摊平申请/建立/归还总线的开销；\kw{交替访存}把一个 CPU 工作周期拆成 CPU 与 DMA 两个分周期，CPU 几乎不等待但要求 CPU 工作周期长于主存存取周期，硬件最复杂。三者不是抽象的优劣排序，必须根据题面给出的设备周期、主存周期和块长度判断成立条件。

\section{通道是 DMA 之上的组织层}
通道不是普通 DMA 控制器的别名。DMA 的行为由 AR/WC/CR 等参数固定，通道则拥有自己的通道指令系统，能从主存取出通道程序，自主完成设备选择、块组织和错误处理。层级可写为
\[
\text{CPU}\to\text{Channel}\to\text{Device Controller}\to\text{Device}.
\]
字节多路通道按字节交叉服务大量低速设备；选择通道按块独占通道服务少数高速设备；数组多路通道按块交替并允许多个子通道，折中吞吐与并行。通道属于 I/O 系统的扩展知识，若题目只考 DMA，仍应把它作为“DMA 已卸掉搬运、通道进一步卸掉组织”的边界，而不把 OS 请求队列写成硬件动作。

\subsection{I/O 方式的统一定量骨架}
若设备产生速率为 $R$ B/s、每次端口/缓冲传送 $u$ B，则程序查询和程序中断的介入频率都约为 $R/u$；查询的 CPU 占比可写成 $(R/u)C_{poll}/f$，中断的占比为 $(R/u)m_{irq}/f$，其中 $m_{irq}$ 包含响应、保护现场、搬运、恢复和返回。DMA 以块大小 $B_s$ 介入，设置和完成的次数约为 $R/B_s$，占比应写成
\[
\rho_{DMA}=\frac{(R/B_s)\,m_{setup+finish}}{f},
\]
并把 Cache 一致性、描述符、仲裁和后处理成本单列。一次块传送通常是 CPU 预处理一次、完成中断/后处理一次，所以“DMA 介入次数与数据个数无关”只在块大小和完成模型已声明时成立。

查询、中断、DMA 和通道可沿“CPU 介入程度”排列：查询全程等待，中断通常每个数据/小单元介入，DMA 每块配置与收尾，通道每个 I/O 任务只需启动和完成交接。升级不是绝对的“越快越好”：低速设备可能以查询的硬件成本最优；高速设备若 $T_{device\ gap}<T_{irq\ response}+T_{ISR}$，中断在物理上会丢数据，应采用 DMA 或更强缓冲。

\begin{methodbox}{陌生总线/I/O 题的停止条件}
完成一题前必须写出：当前 initiator/target、请求对象（CPU 执行权还是总线控制权）、事务四阶段、数据路径、完成证据、异常/中断入口和软件接手位置。若只写“DMA 更快”“中断更高效”而没有给出块大小、端口宽度、仲裁与等待周期，定量结论还未成立。
\end{methodbox}

\subsection{总线标准与并行/串行演进}
总线标准不是只规定带宽，而是同时规定四类接口契约：机械特性（插头、插槽、引脚几何）、电气特性（方向、电平和信号完整性）、功能特性（每根线承担地址/数据/控制哪种语义）和时间特性（何时有效、何时采样）。缺少任何一层，模块都可能“插得上却不能通信”。

低频短距离下，并行线一次传多位；高频长距离下，走线长度差、串扰和时钟偏移会使最快/最慢位无法在同一采样窗内稳定到达，频率越高问题越严重。串行差分链路以较少的位间对齐约束支持更高频率，并可从数据跳变恢复时钟；多条独立 lane 再把吞吐量线性叠加。因此“并行一定比串行快”不是稳定结论，必须连同距离、频率和信号完整性条件判断。

从共享并行总线转向点对点串行链路，还减少了多设备负载、共享仲裁和带宽分摊。PCIe 的 $\times n$ 表示 $n$ 条 lane，而不是一条 $n$ 位并行数据总线；USB、SATA 等也体现串行链路的工程取舍。具体版本、编码比例和标准速率属于题设实例，不应升级为 408 的固定常数，但“并行→串行、共享→点对点”的因果链应保留。

\end{document}
